\documentclass[a4paper,12pt]{article}
\usepackage[T1]{fontenc}
\usepackage[utf8]{inputenc}
\usepackage[polish]{babel}
\usepackage{geometry}
\usepackage{hyperref}
\usepackage{enumitem}

\geometry{margin=2.5cm}
\hypersetup{colorlinks=true, linkcolor=black, urlcolor=blue}

\raggedbottom
\setlength{\parskip}{0.35em}
\setlength{\parindent}{0pt}

\newcommand{\nextitem}{\par\vspace{0.85cm}}

\newcommand{\protokolMeta}[2]{%
  \begin{center}
    {\Large\bfseries PROTOKÓŁ ZE SPOTKANIA SCRUMOWEGO}\\[0.35em]
  \end{center}
  \vspace{0.5cm}
  \textbf{Data:} #1\par\vspace{0.35cm}
  \textbf{Uczestnicy:} #2\par
}


\begin{document}

\protokolMeta{2026-10-21}{Janusz Fijałkowski, Konrad Fijałkowski, Anna Koncewicz, Jakub Grędzicki}

\nextitem
\textbf{1. Co zostało zrealizowane od ostatniego spotkania:}
\par\vspace{0.3cm}
Utworzenie kompletu user stories --- usunięcie duplikatów, doprecyzowanie.

Szablon frontendu w repo (\#33 zamknięte).

Wycena pierwszych wymagań.

Nowe / doprecyzowane issues tego dnia m.in.\ \#37--\#45 (profil/zdjęcia, tokeny, \textbf{grzybobrania} \#43 + split FE \#44 / BE \#45, członkowie \#40, zadania \#42).

\nextitem
\textbf{2. Nad czym obecnie pracujemy:}
\par\vspace{0.3cm}
Kryteria akceptacji do wymagań.

Umieszczenie wymagań w kontekście domenowym (grzybobrania = projekty lab.).

Wymagania konieczne do MVP (rdzeń: CRUD grzybobrań + zadania + członkowie).

DoR: wycena i priorytety.

\nextitem
\textbf{3. Napotkane przeszkody:}
\par\vspace{0.3cm}
Brak zgłoszonych przeszkód.

\nextitem
\textbf{4. Data kolejnego spotkania:}
\par\vspace{0.3cm}
2026-10-23, MS Teams

\end{document}
